% !TeX program = lualatex
% Compile from the repository root: lualatex --interaction=nonstopmode --halt-on-error --synctex=1 --output-directory=data data/storm-table-check.tex
% Compiling from data/ or with pdfLaTeX also works.
% Thesis layout, as recorded in storm-rew-bounded-benchmarking/test-results.tex.
\documentclass[fontsize=10pt,paper=b5,twoside,DIV=11,BCOR=9mm,
    headinclude=true,footinclude=false]{scrartcl}
\usepackage{iftex}
\ifPDFTeX
  \usepackage[T1]{fontenc}
  \usepackage{libertine}
\else
  \usepackage[no-math]{fontspec}
  \setmainfont{Linux Libertine O}
  \setsansfont{Linux Biolinum O}
  \setmonofont[Scale=0.875]{Droid Sans Mono}
\fi
\usepackage[libertine]{newtxmath}

% Resolve data files from either the repository root or the data/ directory.
\makeatletter
\def\input@path{{data/}}
\makeatother
\usepackage[table]{xcolor}
\usepackage{adjustbox}
\usepackage{array}
\usepackage{booktabs}
\usepackage{ifthen}
\usepackage{multirow}
\usepackage{pgfplots}
\pgfplotsset{compat=1.18}
\usepgfplotslibrary{statistics}
% RWTH color scheme
% see: http://www9.rwth-aachen.de/global/show_document.asp?id=aaaaaaaaaadrbpl#page=27

\newcommand{\CMYKCOLORS}{}
\ifdefined\CMYKCOLORS
% RWTH blue
\definecolor{RWTHblue}{cmyk}{1,0.5,0,0}
\definecolor{RWTHlightblue}{cmyk}{0.45,0.14,0,0} % RWTH blue 50%

% Main colors
\definecolor{RWTHblack}{cmyk}{0,0,0,1}
\definecolor{RWTHmagenta}{cmyk}{0,1,0.25,0}
\definecolor{RWTHyellow}{cmyk}{0,0,1,0}

\definecolor{RWTHpetrol}{cmyk}{1,0.3,0.5,0.3}
\definecolor{RWTHteal}{cmyk}{1,0,0.4,0}
\definecolor{RWTHgreen}{cmyk}{0.7,0,1,0}
\definecolor{RWTHmaygreen}{cmyk}{0.35,0,1,0}
\definecolor{RWTHorange}{cmyk}{0,0.4,1,0}
\definecolor{RWTHred}{cmyk}{0.15,1,1,0}
\definecolor{RWTHbordeaux}{cmyk}{0.25,1,0.7,0.2}
\definecolor{RWTHviolet}{cmyk}{0.7,1,0.35,0.15}
\definecolor{RWTHpurple}{cmyk}{0.6,0.6,0,0}

% Other colors
\definecolor{RWTHblue75}{cmyk}{0.75,0.38,0,0} % RWTH blue 75%
\definecolor{RWTHblue25}{cmyk}{0.23,0.07,0,0} % RWTH blue 25%
\definecolor{RWTHsuperlightblue}{cmyk}{0.09,0.03,0,0} % RWTH blue 10%
\definecolor{RWTHwhite}{cmyk}{0,0,0,0}
\definecolor{RWTHgrey}{cmyk}{0,0,0,0.75}
\definecolor{RWTHlightgrey}{cmyk}{0,0,0,0.25}
\definecolor{RWTHsuperlightgrey}{cmyk}{0,0,0,0.01}
\else

% RWTH blue
\definecolor{RWTHblue}{RGB}{0,84,159}
\definecolor{RWTHlightblue}{RGB}{142,186,229} % RWTH blue 50%

% Main colors
\definecolor{RWTHblack}{RGB}{0,0,0}
\definecolor{RWTHmagenta}{RGB}{227,0,102}
\definecolor{RWTHyellow}{RGB}{255,237,0}

\definecolor{RWTHpetrol}{RGB}{0,97,101}
\definecolor{RWTHteal}{RGB}{0,152,161}
\definecolor{RWTHgreen}{RGB}{87,171,39}
\definecolor{RWTHmaygreen}{RGB}{189,205,0}
\definecolor{RWTHorange}{RGB}{246,168,0}
\definecolor{RWTHred}{RGB}{204,7,30}
\definecolor{RWTHbordeaux}{RGB}{161,16,53}
\definecolor{RWTHviolet}{RGB}{97,33,88}
\definecolor{RWTHpurple}{RGB}{122,111,172}

% Other colors
\definecolor{RWTHblue75}{RGB}{64,127,183} % RWTH blue 75%
\definecolor{RWTHblue25}{RGB}{199,221,242} % RWTH blue 25%
\definecolor{RWTHsuperlightblue}{RGB}{232,241,250} % RWTH blue 10%
\definecolor{RWTHwhite}{RGB}{255,255,255}
\definecolor{RWTHgrey}{RGB}{51,51,51}
\definecolor{RWTHlightgrey}{RGB}{204,204,204}
\definecolor{RWTHsuperlightgrey}{RGB}{247,247,247}
\fi

\newcommand{\model}[1]{\textsc{#1}}
\newcommand{\resolution}{\ensuremath{\eta}}
\newcommand{\sizethreshold}{\ensuremath{\tau}}
\newcommand*{\config}[2][]{\ifthenelse{\equal{#1}{}}{\textsc{#2}}{\textsc{#2}$_{#1}$}}

\newcommand{\checktable}[1]{%
  \par\begingroup\centering
    \input{#1}\par
  \endgroup\clearpage
}

% Captions remain in the including document; generated panels embed their data.
\newcommand{\checkplot}[3]{%
  \par\begingroup\centering
    {\footnotesize #2}\par\medskip
    {\small\textbf{#1} #3\par}%
  \endgroup\clearpage
}

% All retained tables, followed by the selected approximation figures.
\setlength{\parindent}{0pt}
\raggedbottom

\begin{document}
\checktable{tables-general/tablebenchmarks.tex}
\checktable{tables-general/tableprbmin.tex}
\checktable{tables-general/tableprbmax.tex}
\checktable{tables-general/tablerewmin.tex}
\checktable{tables-general/tablerewmax.tex}
\checktable{tables-general/overapprox/tableoverapproxprbmin.tex}
\checktable{tables-general/overapprox/tableoverapproxprbmax.tex}
\checktable{tables-general/overapprox/tableoverapproxrewmin.tex}
\checktable{tables-general/overapprox/tableoverapproxrewmax.tex}
\checktable{tables-general/overapprox/tableoverapproxresolutionprbmin.tex}
\checktable{tables-general/overapprox/tableoverapproxresolutionprbmax.tex}
\checktable{tables-general/overapprox/tableoverapproxresolutionrewmin.tex}
\checktable{tables-general/overapprox/tableoverapproxresolutionrewmax.tex}
\checktable{tables-general/underapprox/tableunderapproxthreshold-examples.tex}
\checktable{tables-general/underapprox/tableunderapproxresolution-examples.tex}
\checktable{tables-general/underapprox/tableunderapproxprbmin.tex}
\checktable{tables-general/underapprox/tableunderapproxprbmax.tex}
\checktable{tables-general/underapprox/tableunderapproxrewmin.tex}
\checktable{tables-general/underapprox/tableunderapproxrewmax.tex}
\checkplot{Discretisation: improvement versus runtime.}{%
  \input{tikzpictures/discretisation-budget.tex}%
}{%
  Each curve counts instances for which at least one successful run of that mode,
  across resolutions 2--12, strictly tightens the MDP bound. Values are compared at
  precision $10^{-6}$; each instance is counted once per mode. The horizontal axis
  is the wall-clock time of an individual run, limited to 1,800\,s. Configuration
  selection is retrospective; the cost of searching configurations and computing
  the MDP baseline is excluded.
}
\checkplot{Discretisation: resolution trade-off.}{%
  \input{tikzpictures/discretisation-resolution.tex}%
}{%
  Static/dynamic boxes appear left/right at each resolution. Boxes show quartiles,
  medians, and minimum--maximum whiskers. Top: improvement over MDP normalised by
  each instance's best observed gain across both modes and all resolutions;
  only instances with a successful MDP baseline and positive gain are included.
  Bottom: successful wall-clock runtimes across all instances on a logarithmic
  axis. Runs must finish within 1,800\,s. Counts above the panels show sample
  sizes; failed and unrun cases are omitted from the boxes.
}
\checkplot{Under-approximation: CUT versus CLIP runtime.}{%
  \input{tikzpictures/cut-vs-clip-runtime.tex}%
}{%
  Matched exploration thresholds $\tau=2^{12}$, with clipping resolution
  $\eta=2$. Circles denote probability objectives and triangles reward objectives.
  Above the solid equality diagonal CLIP takes longer; dotted lines show
  factor-of-two differences. Successful runs finish within 1,800\,s. Times at most
  0.1\,s are grouped at the lower edge for display. The separate n/a edges denote
  failed or unavailable runs and have no numeric runtime meaning. Bounds at
  these settings can differ; this compares cost at equal exploration parameters,
  rather than equal bound quality.
}
\checkplot{Under-approximation: Drone threshold trade-off.}{%
  \input{tikzpictures/drone-threshold.tex}%
}{%
  Drone with $N=5$ and $R=3$, maximising reachability probability. CUT and CLIP
  vary the exploration threshold from $2^8$ through $2^{32}$, showing CLIP at
  resolutions $\eta=2$ and $\eta=3$. Top: probability lower bound (larger is
  tighter). Bottom left: wall-clock runtime. Bottom right: resulting belief-MDP
  state count. Both lower panels use logarithmic vertical axes. Numeric points
  are independent successful runs completing within 1,800\,s and are shown
  without connecting lines. Missing measurements are omitted.
  The separate TO and MO$^*$ rows show categorical outcomes, never numeric
  measurements. MO$^*$ groups allocation failures and non-timeout SIGKILLs,
  following the tool-paper convention; a kill alone does not establish its cause.
  CUT reaches 33.5 million states at $\tau=2^{25}$ in 542\,s;
  at $2^{26}$ it is killed after 546\,s, well before the time limit.
  Across all 950 explicit-threshold CUT runs, 177 are killed and only 2 time out.
}
\checkplot{CUT outcomes across all benchmark instances: heatmap.}{%
  \input{tikzpictures/cut-outcomes-heatmap.tex}%
}{%
  Each row is one benchmark instance and objective, grouped by benchmark family;
  columns show explicit thresholds $2^8$ through $2^{32}$. Cells show success,
  timeout (TO), or MO$^*$. The 1,800\,s limit applies to each run. MO$^*$ groups
  allocation failures and non-timeout SIGKILLs; a kill alone does not establish
  its cause. All 38 instances are included. At $2^{32}$, 25 runs are classified
  MO$^*$ and 13 succeed. Across the sweep, 177 runs are killed and 2 time out.
}
\checkplot{CUT: bound improvement and exploration limits across all instances.}{%
  \input{tikzpictures/cut-bound-overview.tex}%
}{%
  Top: each successful CUT bound $b$ is compared with that instance's tightest
  successful bound $b^\star$ across explicit thresholds $2^8$--$2^{32}$, using
  $100|b-b^\star|/\max(|b|,|b^\star|)$ (zero when both are zero).
  Lower is better; zero matches the best observed bound, rather than certifying
  optimality. Comparisons use precision $10^{-6}$. Boxes show quartiles and
  medians, with minimum--maximum whiskers. All 38 instances are eligible,
  including those unchanged from the smallest threshold. Failed runs supply no
  bound and are counted in the lower panel, so the sample changes with threshold.
  Of the 38 instances, 31 obtain a tighter bound than at $2^8$; all 38 runs
  still succeed at $2^{24}$. The 1,800\,s per-run limit applies throughout.
  MO$^*$ includes non-timeout SIGKILLs, whose cause is not confirmed by the logs.
}
\checkplot{Combining bounds: interval tightness across all instances.}{%
  \input{tikzpictures/combined-bounds.tex}%
}{%
  Counts of all 38 instances whose paired bounds satisfy the relative-width
  tolerance $100(U-L)/\max(|L|,|U|)$; higher is better. Matching bounds count at
  every positive tolerance. Best tested configurations use successful runs
  within 1,800\,s, preferring exact results and breaking ties by runtime;
  configuration-search time is excluded. Each method supplies its designated
  interval side, reversed for minimisation. DISCR includes both modes;
  CUT/CLIP include heuristic thresholds. CUT is used where CLIP fails.
  Bounds are rounded to precision $10^{-6}$; both-zero endpoints have zero width.
}
\checkplot{Discretisation error when CUT determines the optimum.}{%
  \input{tikzpictures/discretisation-exact-reference.tex}%
}{%
  All 13 instances with an exact CUT result $v^\star$, sorted by
  $100|b_{\mathrm{DISCR}}-v^\star|/\max(|b_{\mathrm{DISCR}}|,|v^\star|)$.
  CUT's error is zero throughout; DISCR retains positive error on 9 instances
  and matches on 4. Discretisation accounts for the remaining interval on this
  exact-reference subset. The best successful DISCR configuration within
  1,800\,s is selected across both modes and resolutions 2--12; the exact CUT
  reference must also complete within 1,800\,s. Bounds are rounded to precision
  $10^{-6}$; both-zero values have zero error. The other 25 instances have no
  exact CUT reference and are excluded from this figure.
}
\end{document}
